% Source: https://artofproblemsolving.com/community/c3h2301154_puzzle_fun
% Stable identity: aops-topic-2301154-q01
\textbf{(a)} Take the triangle $T_1$ with vertices $(0,0)$, $(10,0)$, $(5,8)$ and the long thin triangle $T_2$ with vertices $(-1,3)$, $(-1,4)$, $(12,3.5)$, which pierces $T_1$ from left to right.

\begin{asy}
usepackage("amsmath");
size(8cm);
defaultpen(fontsize(10pt));
pair A1=(0,0), A2=(10,0), A3=(5,8);
pair B1=(-1,3), B2=(-1,4), B3=(12,3.5);
fill(A1--A2--A3--cycle,lightgray); fill(B1--B3--B2--cycle,lightgray);
draw(A1--A2--A3--cycle,blue); draw(B1--B3--B2--cycle,red);
pair X1=extension(A1,A3,B1,B3), X2=extension(A2,A3,B1,B3), X3=extension(A2,A3,B2,B3), X4=extension(A1,A3,B2,B3);
draw(A1--A2--X2--B3--X3--A3--X4--B2--B1--X1--cycle,linewidth(1.2));
dot(A1);dot(A2);dot(A3);dot(B1);dot(B2);dot(B3);dot(X1);dot(X2);dot(X3);dot(X4);
label("$X_1$",X1,SE,fontsize(7pt)); label("$X_2$",X2,SW,fontsize(7pt)); label("$X_3$",X3,1.6*dir(151),fontsize(7pt)); label("$X_4$",X4,NE,fontsize(7pt));
label("$(0,0)$",A1,SW,fontsize(8pt)); label("$(10,0)$",A2,SE,fontsize(8pt)); label("$(5,8)$",A3,N,fontsize(8pt));
label("$(-1,3)$",B1,W,fontsize(8pt)); label("$(-1,4)$",B2,W,fontsize(8pt)); label("$(12,3.5)$",B3,E,fontsize(8pt));
\end{asy}

The two long sides of $T_2$ (on the lines $y=3+\dfrac{x+1}{26}$ and $y=4-\dfrac{x+1}{26}$) each cross the left side $y=1.6x$ and the right side $y=1.6(10-x)$ of $T_1$, at points with heights between $3$ and $4$; the base and the apex of $T_1$ lie outside $T_2$, and all three vertices of $T_2$ lie outside $T_1$. Going around the union, the outline passes through
\[(0,0),\ (10,0),\ X_2,\ (12,3.5),\ X_3,\ (5,8),\ X_4,\ (-1,4),\ (-1,3),\ X_1,\]
where $X_1,X_2,X_3,X_4$ are the four crossing points. No two consecutive sides are collinear, so the outline is a polygon with $\boxed{10}$ sides (shown thick).

\textbf{(b)} Every vertex of the outline is a point where the outline turns. Such a point lies on the boundary of $T_1$ or of $T_2$. If it is not a vertex of $T_1$ or of $T_2$, then the outline turns there from a side of one triangle to a side of the other, so it is a crossing point of a side of $T_1$ with a side of $T_2$ that are not on the same line.

Count these crossing points on one side $s$ of $T_1$. The part of $s$ inside $T_2$ is a single segment (or empty), because $T_2$ is convex. At a crossing point that is not a vertex of $T_2$, the side $s$ goes from inside $T_2$ to outside $T_2$, so the point is an endpoint of that segment. Hence $s$ contains at most $2$ such points, and the three sides of $T_1$ contain at most $6$.

So the outline has at most $3$ vertices of $T_1$, $3$ vertices of $T_2$ and $6$ crossing points, that is, at most $12$ vertices, and therefore at most $12$ sides. $\blacksquare$ (Twelve is possible: two equilateral triangles with the same centre, one rotated by $180^\circ$, form a six-pointed star with $12$ sides.)
